\documentclass{article}
\usepackage[T1]{fontenc}
\usepackage{lmodern}
\usepackage{graphicx}
\usepackage{amsmath,amssymb}
\usepackage[a4paper,margin=2cm]{geometry}
\usepackage[absolute,overlay]{textpos}
\setlength{\TPHorizModule}{1mm}
\setlength{\TPVertModule}{1mm}
\usepackage[indonesian]{babel}
\usepackage{hyperref}
\setlength{\emergencystretch}{2em}
% unit: Halaman 32

\begin{document}

% === Halaman 32 (llm) ===

\section*{Efek longsoran (avalanche effect)}

\begin{itemize}
    \item Efek longsoran adalah properti kunci yang diinginkan dari algoritma enkripsi,
    \item Tujuan: menunjukkan perubahan satu bit pada plainteks atau kunci menghasilkan perubahan besar pada cipherteks
    \item Efek longsoran membuat upaya untuk mencoba menebak kunci menjadi tidak mungkin
    \item DES menunjukkan efek longsoran yang kuat
\end{itemize}

\end{document}
